\section{Introduction}
A renewal contract may promise service at many public histories while its terminal interface implements only a small number of commands. The provider must choose the commands, allocate the promised service across histories, and pay for any intermediate service that cannot be delivered by the terminal command. Individual expectation caps and realization ceilings make these choices interdependent. A command useful at one history may be inadmissible at another; opening a command changes the common interface rather than a single local allocation. The design problem is therefore joint even when allocating service for a fixed interface is easy.

We study this problem on a finite candidate catalog with a fixed installed-symbol budget and nonnegative opening charges. Our principal structural result is terminal-first pooling. Histories admitting the same catalog prefix can be optimized together without averaging their caps or costs. The exact service-cost correction depends on the selected book only through its last eligible level. An ordered path can consequently charge the correction once and jointly optimize the book and the individual targets. When all histories have common catalog eligibility, this yields an exact polynomial algorithm with variable history count, catalog size, and symbol budget. Common eligibility does not require that every catalog level be eligible.

Catalog eligibility is nevertheless not an intrinsic dimension. Adding an unused candidate can split a class without changing the optimum. We address that dependence in two ways. First, we characterize refinement and threshold stability, provide certified deletion rules, and solve threshold-robust design by its lower-ceiling representation. Second, we develop a selected-prefix search. Its state commits to installed symbols rather than a partition of all histories. A mandatory-prefix extension of the two-sided price graph supplies original-budget bounds, and a free-completion bound permits simpler certificates. The resulting binary cover is exact for every fixed symbol budget, regardless of the number of catalog eligibility classes. Fixed-budget polynomial enumeration was already established in the antecedent complexity proposition; the new contribution is the constrained price recurrence and independently checkable pruning, not that elementary enumeration implication. This is a complementary polynomial regime, not a claim of fixed-parameter tractability when the symbol budget itself grows.

The algorithms retain their original-space interpretation at every interruption. A lower policy satisfies the exact aggregate promise, each individual cap, and every supportwise ceiling. Upper certificates combine scalar support prices, Bellman inequalities, and complete decision-tree coverage. Independent verification does not trust an optimizer status or substitute a larger alphabet. We also retain the group-mean box method as a distinct accuracy-based route and give its explicit dyadic cover bound, including initial widths and rounding.

Our experiments separate arithmetic scale from eligibility geometry. Strict interior prefixes, heterogeneous thresholds within a gap, near-boundary perturbations, and successive catalog insertions test the structural claims. We revisit every distinct difficult instance from the preceding study without replacing its inputs or historical failures. Matched-class experiments vary the proposed structural parameter on otherwise fixed primitives. Direct mixed-integer formulations expose solver times, nodes, numerical gaps, and envelope allowances alongside exact rational certificates. All inputs are declared synthetic. A capacity-reservation model derives the same optimization structure from expected served demand rather than asserting empirical calibration.

The use of capped resource allocation and branch-and-bound is classical. We distinguish these ingredients from the exact last-eligible-level correction, the mandatory-prefix price decomposition, and the resulting original-budget certificate. Existing work on separable resource allocation supplies the inner multiplier machinery \citep{PatrikssonStromberg2015,SchootUiterkamp2022}; branch-and-bound supplies the general partition-and-bound principle \citep{LawlerWood1966}. The substantive question here is which exact bounds and primal reconstructions remain valid when command selection, safety eligibility, heterogeneous service costs, and the promise equality are optimized together.

\paragraph{Relation to quantization and constrained paths.}
Ordered interval optimization and scalar quantization already have dynamic-programming and matrix-searching formulations \citep{Wu1991,NielsenNock2014}. Random splitting between grid points preserves a prescribed mean in dual quantization \citep{PagesWilbertz2012}. Here an adjacent lottery is the canonical local implementation, not the new computational claim. A shared charged catalog, individual realization ceilings, heterogeneous intermediate costs, and an aggregate contractual equality must still be optimized jointly. Our correction retains those individual restrictions while localizing the effect of pooled service at one selected level.

Lagrangian pricing of constrained paths is also established \citep{HandlerZang1980}. We do not infer primal exactness from a price optimum or transfer a generic constrained-path hardness result to this restricted model without a reduction. Instead, the mandatory-prefix graph proves an exact support formula for every partial book, and the binary cover supplies the missing primal completeness. Continuous design, Monge acceleration, guarded aggregation, and harmonic coarsening remain complementary results under their original assumptions; the companion's reading map preserves their complete derivations rather than recasting them as new ingredients of the selected-prefix algorithm.
